\section{一座无限的山：问题不可避免，但问题可以解决}

本章从《The Beginning of Infinity》开始。杨植麟说，问题是不可避免的，但问题是可以解决的；也许这座雪山没有尽头，但他希望它一直没有尽头。这个比喻很适合模型公司：每次模型能力提高，都会解锁新场景，同时展开新问题。预训练、RLHF、长思考、Agentic 能力、工具使用和评价体系，都像雪山不同高度的路线。

\begin{figure}[H]
\centering
\includegraphics[width=0.92\textwidth]{infinite-mountain-map.png}
\caption{一座无限的山：问题不可避免，但问题可以解决；模型进步不断打开新问题。自制概念图，依据 00:00:02--00:05:18 对谈内容整理。}
\end{figure}

\begin{knowledgebox}{读图：无限不是悲观，而是研究氧气}
如果山有尽头，模型公司很快会变成纯工程公司；如果山没有尽头，研究、产品和组织都会持续面对新问题。无限的山意味着长期问题空间仍然存在。
\end{knowledgebox}

\subsection{Test-time scaling：长思考和 Agent 指向同一方向}

上一节讲“问题会继续展开”，本节进入第一个技术判断。杨植麟认为，不管是基于长思考的强化学习，还是 Agent 的强化学习，都指向 test-time scaling。所谓 test-time scaling，就是在推理/测试阶段投入更多计算：更长推理链、更多工具调用、更多验证和更多迭代。

\begin{figure}[H]
\centering
\includegraphics[width=0.96\textwidth]{test-time-scaling-map.png}
\caption{Test-time Scaling：长思考 RL 和 Agent RL 都在扩展推理时计算。自制概念图，依据 00:03:38--00:24:58 对谈内容整理。}
\end{figure}

\begin{knowledgebox}{读图：测试时扩展不是只多输出 token}
长思考让模型在内部推理上花更多 token；Agent 让模型通过工具和环境获得外部反馈。二者都在测试阶段花更多计算，只是一个偏内部思考，一个偏外部行动。
\end{knowledgebox}

\begin{importantbox}{Test-time scaling 的直觉公式}
\[
\text{Capability} \approx f(\text{model}, \text{test-time compute}, \text{feedback})
\]
其中，\(\text{model}\) 是基座能力，\(\text{test-time compute}\) 是推理阶段额外投入，\(\text{feedback}\) 是验证、工具返回或环境反馈。Agentic LLM 的关键是把后两者做成系统。
\end{importantbox}

\subsection{L1 到 L5 不一定串行}

上一节把 test-time scaling 拆成内部思考和外部行动，本节进一步说明这些能力不一定线性爬楼。访谈提到 L1 到 L5 不一定是串行关系，Claude 就是在 reasoning 不一定最强的情况下，在 Agent 上做得很好。这个判断很重要：模型能力发展不是固定阶梯，Reasoning、Agent、工具、产品和反馈可以并行押注。只有当模型参与到开发过程，才可能解锁真正的 Innovator 阶段。

\begin{figure}[H]
\centering
\includegraphics[width=0.96\textwidth]{l1-l5-agent-ladder.png}
\caption{L1 到 L5 不一定串行：Reasoning 与 Agent 可以并行押注，Claude 是典型例子。自制概念图，依据 00:12:00--00:24:58 对谈内容整理。}
\end{figure}

\begin{warningbox}{不要把能力等级当作机械阶梯}
如果把 L1 到 L5 理解为必须逐级通关，就会低估产品和系统路线。Agent 能力可能在某些任务上先跑出来，再反过来推动模型训练和工具生态。
\end{warningbox}

\begin{knowledgebox}{L 阶段的教学解释}
\begin{tabularx}{\textwidth}{p{0.16\textwidth}p{0.34\textwidth}X}
\toprule
阶段 & 直觉 & 关键变量 \\
\midrule
L1/L2 & 聊天、问答、基础 reasoning & 模型底座与对齐。 \\
L3 & 能使用工具完成明确任务 & 工具接口、任务定义、反馈。 \\
L4 & 参与开发、创新和复杂问题解决 & 长程 Agent、验证、代码能力。 \\
L5 & 更高自主性系统 & 安全、组织、权限和持续反馈。 \\
\bottomrule
\end{tabularx}
\end{knowledgebox}

\subsection{本章小结}

“无限的山”给这期提供了精神底色：模型能力每前进一步都会展开新问题；test-time scaling 则给出技术方向：让模型在推理时思考更多、行动更多、获得更多反馈。

